\section{Cohomological operators and deformation}\label{sec:cohomology}

Proposition~\ref{prop:realization} specifies Euler pairings and bounds an
\(\Ext^2\) space.  We now express both conclusions in linear algebra.
The Euler pairings determine an algebraic component of a Mukai vector;
the bound constrains the kernel of the Hochschild cap action on the
entire vector.  When that kernel has the expected dimension, the
semiregularity map is injective, and the complex formally deforms along
its Hodge locus.  Spreading and specialization then transport the
algebraicity of its Chern character.  This is the mechanism used in
Sections~\ref{sec:exactks} and~\ref{sec:auxiliary}.

\subsection{The linear dictionary}

Two gradings will enter the argument.  To choose the symplectic input
class, we use the Lefschetz action on the cohomology of \(D\).  To study
the resulting complex, we use the actual Hodge grading on the mirror.
We first describe the former as a linear model for the latter; identifying
their transcendental Hodge structures will require the later rank argument.

We use the Mukai form \(q\) on the total even cohomology of a K3
surface: it is the intersection form on \(H^2\), and
\(q(1,[\mathrm{pt}])=-1\).  Mukai vectors are
\(v(\mathcal E)=\ch(\mathcal E)\sqrt{\td}\).
When using the Hodge circle, we suppress Tate twists; invariance means
type \((p,p)\) separately in each cohomological degree.

Suppose first that \(Y\) is a quartic.  Recall the isometry
\(\rho:P_Y\longrightarrow P_B\) supplied by the sphere tests in
Section~\ref{sec:setup}, where
\[
 P_Y=\omega_Y^\perp\subset H^2(Y,\Q),\qquad
 P_B=H^0(B,\Q)\oplus\NS(B)_\Q\oplus H^4(B,\Q).
\]
The orthogonal complement of \(P_Y\) in the Mukai space is
\[
 U=\langle 1,\omega_Y,[\mathrm{pt}]\rangle.
\]
It is a split ternary quadratic space.  Lefschetz theory gives an
\(\mathfrak{sl}_2\)-action with Cartan operator \(2-\deg\) and lowering
operator \(f=\omega_Y\cup-\).  The space \(U\) is its irreducible
representation of highest weight two, while \(P_Y\) is trivial.
Set \(t_+=1\) and \(t_-=-[\mathrm{pt}]\), so
\(q(t_+,t_-)=1\), and choose a weight-zero basis vector \(t_0\) of
\(U\).  We identify \(\Spin(U)\) with \(\SL_2\).
Write \(U_B=T(B)\).  Witt cancellation in the two rational Mukai lattices
identifies \(U\), as a quadratic space, with \(U_B\).
This is a rational quadratic identification only; no assertion about
the transcendental action of mirror symmetry is used.

For the torus write
\[
 W=M_0\otimes H,\qquad \dim M_0=g,\qquad
 H=\Q e_+\oplus\Q e_-,\qquad f(e_+)=e_-.
\]
The vectors \(e_+,e_-\) are the standard \(\SL_2\)-basis, of weights
\(1,-1\).  Write \(a_i,b_i\) for the corresponding basis vectors of
\(W\), \(1\le i\le g\).  The creation--annihilation transform is the
linear isomorphism
\begin{equation}\label{coh:transform}
 \Phi:H^*(T_g,\Q)\longrightarrow\bigwedge W
\end{equation}
that sends \(p_1\cdots p_g\) to \(1\), cup product by \(q_i\) to
exterior multiplication by \(b_i\), and cup product by \(p_i\) to
contraction by \(a_i^*\).  With consistent ordering, the other two
creation--annihilation rules are
\[
 \Phi\,\iota_{q_i^*}\,\Phi^{-1}=\iota_{b_i^*},
 \qquad
 \Phi\,\iota_{p_i^*}\,\Phi^{-1}=a_i\wedge-.
\]
Together these four rules and the chosen vacuum characterize the
transform.  Thus the Cartan
operator on the source is \(g-\deg\), and cup product by
\(\omega_T\) becomes \(f\).  The transform is not a ring isomorphism.

For a symmetric integral slope matrix \(M\), the graph class is, up
to a fixed orientation sign, \(\prod_i(p_i-(Mq)_i)\).  Consequently
\begin{equation}\label{coh:graph-exponential}
 \Phi\bigl(\PD[G_M]\bigr)
 =\pm\exp\left(\sum_{i,j}M_{ij}a_i b_j\right).
\end{equation}
One may verify the identity by the annihilating operators of the graph
and its constant coefficient.  Mark \(H^1(E,\Q)\) by \(H\), and hence
\(H^1(E^g,\Q)\) by \(W\).  The numerical graph calculation in
Section~\ref{sec:setup} makes the right side the Mukai vector of the
mirror line, with the common shift understood.  The marking is oriented
or reverse oriented according to the single sign convention for all
graph lines; it is fixed throughout the construction.

The degree-two \(\SL_2\)-invariants in \(\bigwedge W\) are precisely
the expressions \(\sum M_{ij}a_i b_j\) with \(M\) symmetric.  The
first fundamental theorem for \(\SL_2\) says that they generate its
exterior-algebra invariants \cite{WeylClassicalGroups}.  Expanding the exponentials and polarizing
therefore shows that the classes in
\eqref{coh:graph-exponential} span those invariants.  A sufficiently
large finite graph list has the same span.  Since \(E\) is non-CM,
these are exactly its power's rational Hodge classes.  Their Euler
pairing is nondegenerate and agrees, up to one fixed nonzero scalar,
with the intersection pairing under \(\Phi\): on pairs of graph
exponentials this is the determinant formula, and those classes span.

Combine the torus transform with the sphere isometry to define
\begin{equation}\label{coh:dictionary}
 \Psi=(\rho\oplus\id_U)\otimes\Phi:
 H^*(Y,\Q)\otimes H^*(T_g,\Q)
 \xrightarrow{\sim}(P_B\oplus U)\otimes\bigwedge W.
\end{equation}
Here the target still contains the formal Lefschetz space \(U\),
rather than \(U_B\) with its actual Hodge structure.  When \(Y\) is
a point, set \(\Psi=\Phi\).

We can now state precisely what the Euler tests recover.  In the
quartic case let
\begin{equation}\label{coh:input-space}
 C_*\in (P_B\oplus U)\otimes\bigwedge W
\end{equation}
be rational and invariant under the simultaneous \(\SL_2\), acting
trivially on \(P_B\).  In the torus-only case take
\(C_*\in\bigwedge W\) invariant under \(\SL_2\).
Set \(a=\Psi^{-1}(C_*)\).  This lies in \(H^n(D,\Q)\): total Cartan weight
zero means middle degree, and \(\SL_2\)-invariance gives
\(\omega a=0\).
If the annihilator of \(a\) satisfies
\eqref{setup:wedge-injective}, Proposition~\ref{prop:realization}
gives a perfect complex.  Divide its Mukai vector by the common
nonzero rational scalar in the resulting Euler comparison, including
the fixed pairing normalization of \(\Psi\), and call the result \(C\).
Then
\begin{equation}\label{coh:determined-component}
 C|_{P_B}=C_*|_{P_B};
 \qquad\text{in the torus-only case, }C=C_*.
\end{equation}
Indeed the sphere tests lie in \(P_Y\), and the graph tests span all
the relevant torus Hodge classes.  The \(P_B\)-component of an
algebraic class is itself Hodge.  The nondegenerate pairings therefore
determine that component.  They do not yet determine the component in
\(T(B)\).

\subsection{The annihilator and cap operators}

The degree-two annihilator of \(a\) can be calculated directly on
\(C_*\).  A K3 degree-two class becomes
\(x\in P_B\oplus\langle t_0\rangle\).  Its cup action is
\begin{equation}\label{coh:lowering-k3}
 D_x(t_+)=x,\qquad
 D_x(y)=-q(x,y)t_-\quad
   (y\in P_B\oplus\langle t_0\rangle),\qquad
 D_x(t_-)=0.
\end{equation}
The torus degree-two actions are all operators
\begin{equation}\label{coh:lowering-torus}
 F(\lambda,L,\pi)=\lambda\wedge-+F_L+\iota_\pi,
 \quad
 \lambda\in\bigwedge^2\langle b_i\rangle,
 \quad \pi\in\bigwedge^2\langle a_i^*\rangle,
\end{equation}
where \(F_L\) is the derivation sending the \(a\)-copy of \(M_0\)
to the \(b\)-copy by an arbitrary matrix \(L\), and vanishing on the
\(b\)-copy.  On the tensor product the full operator is
\(D_x\otimes1+1\otimes F(\lambda,L,\pi)\).  These formulas follow
from cup product on a surface and the creation--annihilation rules.

The actual Hochschild cap action on \(C\) has the same operator
formulas, after choosing bases for its Hodge weights.  The two
calculations have different inputs:
\[
\begin{array}{c|c|c}
 &\text{formal calculation}&\text{actual calculation}\\ \hline
 \text{class}&C_*&C\\
 \text{K3 complement}&U&U_B=T(B)\\
 \text{weights}&n-\deg\text{ on }H^*(D)&p-q\text{ on }H^{p,q}\\
 \text{operators}&H^2(D)\text{ by cup product}
   &\HH^2(B\times E^g)\text{ by cap product}
\end{array}
\]
Omit the K3 row and the factor \(B\) in the torus-only case.
Equality of these operator formulas does not yet identify \(C\) with
\(C_*\), or the Hodge circle of \(U_B\) with the formal action on
\(U\).  It will let us constrain the unknown component of \(C\)
using the small \(\Ext^2\) space.  We now justify the actual formulas
directly on the mirror, without a Hochschild comparison under HMS.

Use Todd-modified Hochschild--Kostant--Rosenberg isomorphisms.  On
homology the modification wedges with \(\sqrt{\td}\); on cohomology
use the corresponding contraction modification.  The module theorem
of Calaque--Rossi--Van den Bergh identifies cap product with the
polyvector action on Hodge cohomology
\cite[Theorem~1.4]{CalaqueRossiVanDenBergh}.  It sends the Hochschild
class of a perfect complex to its Mukai vector.  We take normalized HKR,
with the factor \(1/p!\) in form degree \(p\), and totalized form
signs: holomorphic forms have Hochschild homological degree one;
antiholomorphic cochains precede polyvector contractions, with the
Koszul signs for moving each contraction through a cochain.

For this actual calculation, reuse the weight notation by choosing
\(t_+\in H^{2,0}(B)\) and
\(t_-\in H^{0,2}(B)\) with \(q(t_+,t_-)=1\).
Choose \(t_0\) in the remaining Hodge line of \(U_B\).
These are complex Hodge directions, in place of the formal basis
\(1,\omega_Y,-[\mathrm{pt}]\) used above.  The weight-zero space is
now the full \((p,p)\)-part of the Mukai space.  The three summands
\[
 H^2(B,\cO_B),\qquad H^1(B,T_B),\qquad
 H^0(B,\bigwedge^2 T_B)
\]
act as \eqref{coh:lowering-k3}, with \(x\) respectively in
\(H^4\), \(H^{1,1}\), and \(H^0\).  Here \(T_B\) in the displayed
groups is the tangent bundle.  Wedge product gives the first action;
contraction with a tangent-vector-valued one-form gives the middle
derivation, skew for the Mukai form; bivector contraction gives the
last action.  These observations verify the relative signs in
\eqref{coh:lowering-k3}, up to an immaterial rescaling of each acting
coordinate.  The same form calculus on the abelian factor gives
\eqref{coh:lowering-torus}, now with \(a,b\) the Hodge directions of
weights \(1,-1\).

There are no cross terms involving the K3 factor, because its
\(\HH^1\) vanishes.  Consequently
\begin{equation}\label{coh:hh-dimension}
 \dim\HH^2(B\times E^g)=22+\binom{2g}{2}=b_2(Y\times T_g),
\end{equation}
with the term \(22\) omitted in the torus-only case.  Passing from
Dolbeault-normalized to Betti-normalized classes does not alter these
rank calculations.  In an output component of type \((j,j+2)\), an
acting polyvector of degree \(b\in\{0,1,2\}\) uses the Mukai
component of half-degree \(j+b\).  The powers of the comparison unit
\(\pm2\pi i\) thus amount to invertible row and column rescalings.

\begin{lemma}\label{coh:rank-trace}
Let \(\mathcal E\) be supplied by
Proposition~\ref{prop:realization}, and let \(C\) be its normalized
Mukai vector.  The kernel of the operators
\eqref{coh:lowering-k3}--\eqref{coh:lowering-torus} acting on \(C\)
has dimension at least \(\dim\mathcal Z\).  If it has dimension
exactly \(\dim\mathcal Z\), the semiregularity trace on
\(\Ext^2(\mathcal E,\mathcal E)\) is injective.
\end{lemma}

\begin{proof}
The cap action on the class of an object factors as
\begin{equation}\label{coh:trace-factorization}
 \HH^2\longrightarrow\Ext^2(\mathcal E,\mathcal E)
       \longrightarrow\HH_{-2}.
\end{equation}
The first map is evaluation of a Hochschild operator on the object;
the second is the boundary--bulk trace.  Under HKR this trace is the
collection of traces against powers of the Atiyah class, followed by
an invertible Todd modification.  These are the diagonal-adjunction
and Atiyah--Chern character identities
\cite{CaldararuMukaiII,BuchweitzFlennerHH}.
They apply to all the components of the trace, not only to
\(H^2(\cO)\).  Rescaling the Mukai vector has no effect on rank.

By \eqref{coh:hh-dimension} and Proposition~\ref{prop:realization},
the composite in \eqref{coh:trace-factorization} has rank at most
\(b_2(D)-\dim\mathcal Z\).  This proves the kernel bound.  If
equality holds, the composite has rank equal to the upper bound on
the dimension of its middle space.  Both ranks must then equal that
dimension, so the trace out of the middle space is injective.
\end{proof}

\subsection{Deforming a complex along its Hodge locus}

We use the following precise form of semiregularity.  It isolates the
only deformation-theoretic input needed after the kernel calculations.
The obstruction principle is due to Buchweitz--Flenner; the version
for all infinitesimal obstructions and varying ambient varieties is
also a consequence of Pridham's theorem
\cite{BuchweitzFlennerSemiregularity,PridhamSemiregularity}.

\begin{proposition}\label{prop:semiregularity}
Let \(\mathscr X\to S\) be a smooth projective family over a smooth
complex algebraic base, let \(s\in S\), and let
\(\mathcal E\in\Perf(\mathscr X_s)\).  Suppose the parallel transport
of \(\ch(\mathcal E)\) remains Hodge on the germ of \(S\) at \(s\),
and that the map
\[
 \Ext^2(\mathcal E,\mathcal E)
 \longrightarrow\bigoplus_{j\ge0}H^{j+2}(\mathscr X_s,\Omega^j),
 \qquad o\longmapsto
  \bigl(\Tr(o\,\operatorname{At}(\mathcal E)^j)/j!\bigr)_j
\]
is injective.  Then \(\mathcal E\) extends to a compatible family of
perfect complexes over the formal germ of \(S\) at \(s\).

If \(S\) is connected and the Chern character is a global parallel
section remaining Hodge on \(S\), that section is algebraic on every
fiber.  The same conclusions hold for the Mukai vector in a family
whose Todd classes are parallel algebraic classes.
\end{proposition}

\begin{proof}
We first explain the obstruction calculation.  Resolve the initial
perfect complex by a bounded complex of vector bundles, and use the
Dolbeault resolution of its graded smooth bundle.  A deformation is
represented by a total-degree-one \((0,*)\)-superconnection squaring
to zero.  Relative smooth forms may be used over the Artin
neighborhoods of \(s\): in a smooth local trivialization take Taylor
coefficients in base variables and set antiholomorphic base variables
to zero.  Local holomorphic coordinates identify these complexes with
relative Dolbeault and de Rham resolutions.  Smooth projectivity gives
Hodge degeneration and cohomology base change.

Consider a small extension of Artin neighborhoods, with kernel \(I\)
annihilated by the maximal ideal.  Lift an already constructed
superconnection arbitrarily, retaining the relative Dolbeault Leibniz
rule.  Its square \(o\) is an \(I\)-valued cocycle in the central
endomorphism complex and represents the obstruction.  Add a smooth
\((1,0)\)-connection and denote the curvature of the resulting full
superconnection by \(R\).  The closed forms \(\str(R^m)\) represent
the horizontal Chern character, up to the usual nonzero normalization;
this follows from the superconnection Chern--Weil transgression
formula.

Sort the curvature by holomorphic form degree.  Its degree-zero part
is \(o\), and the other parts have degrees one and two.  Since
\(I^2=0\), every term with two copies of \(o\) vanishes.  Hence
\(\str(R^m)\) has no component below holomorphic degree \(m-1\),
and its component of degree \(m-1\) is
\begin{equation}\label{coh:obstruction-trace}
 m\,\str\bigl(o(R^1_0)^{m-1}\bigr),
\end{equation}
where \(R^1_0\) is the degree-one curvature at the central fiber.
It represents the Atiyah class.  The horizontal Chern character lies
in filtration \(F^m\) by hypothesis.  Hodge degeneration therefore
makes \eqref{coh:obstruction-trace} zero in Dolbeault cohomology.
This remains true with coefficients in \(I\) on the central fiber,
by freeness and cohomology base change.  All semiregularity traces of
the obstruction vanish.  Injectivity makes \(o\) exact, so a
correction of the lifted superconnection makes its square zero.
Iterating over the powers of the maximal ideal gives compatible lifts.

These lifts are perfect complexes, not merely formal differential
operators.  Locally on a coordinate disk, assume inductively that the
operator modulo \(I\) has been gauged to ordinary Dolbeault plus a
holomorphic differential.  Lift that gauge and an arbitrary
holomorphic matrix lift of the differential.  The difference has
coefficients in \(I\).  Its highest positive form-degree component is
Dolbeault closed by the square-zero equation.  A local primitive gives
a gauge transformation removing it.  Continue downward in form
degree.  The remaining square-zero operator is ordinary Dolbeault
plus a holomorphic differential.  Thus it is locally the Dolbeault
resolution of a perfect complex.  The relative forms are flat over
the Artin base, so the chosen reductions compute derived restriction.
GAGA identifies these analytic lifts with algebraic perfect complexes.

The fixed bounded graded bundle in this local straightening gives
a uniform finite Tor amplitude for the compatible lifts.
Grothendieck existence algebraizes the resulting pseudocoherent tower
over the completed local base
\cite[Theorem~38.29.6, Tag~0DIG]{PerfectGrothendieckExistence}.
The family is proper, flat, and finitely presented, and the restriction
to the closed fiber is perfect.  The relative-perfectness criterion
\cite[Lemma~37.71.3, Tag~0DJZ]{PerfectGrothendieckExistence} therefore
makes the algebraization perfect relative to the completed base.
That base is regular and the family is smooth, so the total space is
regular; the resulting bounded coherent complex is thus perfect.
Spread its finite presentation to a finitely generated subalgebra of
the completion over an affine neighborhood of \(s\), retaining the
closed point and its original complex.  The resulting finite-type
parameter is dominant over that neighborhood: the base is smooth at
\(s\), and its local ring maps injectively to its completion.
On the relevant irreducible component, topological constancy of Chern
characters identifies the classes of the spreading complex with the
pullback of the prescribed parallel section.  Its image contains a
dense open subset of \(S\).

For completeness, the last passage to every fiber uses proper cycle
parameters, not constancy of a chosen complex.  Express algebraicity
of the fixed rational parallel class by finite sums and differences
of cycles with a common positive denominator.  Relative projective
Hilbert schemes give countably many proper parameter spaces for such
expressions; restrict to connected components on which the cycle
class is the prescribed one.  Their images in \(S\) are closed.
A dense open subset of an irreducible complex variety cannot be
covered by countably many proper closed subvarieties.  One such
proper image is therefore all of \(S\).  This proves algebraicity on
every fiber.  A connected smooth base is irreducible componentwise
with no intersecting components, so the stated connected case follows.
Finally, multiplication by \(\sqrt{\td}\) is an invertible parallel
algebraic operation in the indicated families and changes none of
the conclusions.
\end{proof}

The remaining task is now finite-dimensional: choose \(C_*\) so that
the annihilator has the required wedge-injectivity, and show that the
unknown Mukai component is forced by the cap-kernel bound.  Once the
bound is sharp, Proposition~\ref{prop:semiregularity} transports the
resulting algebraic tensor through the desired period family.
